\chapter{Introducción}
En las últimas décadas, la ciencia de materiales ha presentado transformaciones significativas impulsadas especialmente por los avances en nanotecnología. Esta disciplina, capaz ya no solo de analizar, sino también de diseñar y estructurar materiales en escalas invisibles al ojo humano haciendo uso de nuevas y diferentes técnicas, tiene como objetivo configurar propiedades de los mismos de manera controlada. Esto constituye una revolución en un mundo que exige más y más soluciones específicas a problemas complejos en diferentes ámbitos, en los cuales el estudio de nanomateriales tiene un papel importante a desempeñar.

Bajo este contexto, este trabajo se engloba en un tipo específico de nanomaterial, los polímeros nanocompuestos. Estos están construidos a partir de matrices poliméricas sobre las cuales se integra material nanoparticulado. Existen muchos tipos y configuraciones de los mismos, buscando en ellos diferentes características, definidas por las variables usadas en su fabricación y aplicación: el tipo de polímero y nanorelleno utilizado, concentraciones, métodos de unión, entre muchas otras cuestiones que hacen que cada composición varíe sus propiedades respecto a otras y a sus propios componentes.

Con el objetivo de obtener propiedades que sean útiles en muchos campos; desde el ámbito académico, pasando por la industria y el desarrollo militar, hasta soluciones en el campo de la salud y en materia ambiental: constantemente se están planteando, estudiando y manufacturando nuevos tipos de polímeros nanocompuestos.  

Una de las características de mayor interés en este tipo de materiales es la eficiencia de apantallamiento (SE) frente a interferencia electromagnética (EMI). Los dispositivos electrónicos actuales requieren materiales que los protejan del ruido externo y preserven la integridad de las señales que manejan. La proliferación de sistemas que operan en bandas de alta frecuencia ha incrementado los requerimientos de apantallamiento [12], tanto por la mayor facilidad con que los sistemas se acoplan con señales parásitas en ese régimen como por los estándares de compatibilidad electromagnética (EMC). Para aplicaciones prácticas, se considera un umbral de referencia de aproximadamente 20 dB de SE como requisito para aplicaciones comerciales [4]. 

Frente a esto, la tarea de diseñar polímeros nanocompuestos capaces de aislar los sistemas de este tipo de ruido es crucial, pues aporta ventajas estructurales, económicas y de otros tipos, frente a aislamientos convencionales. La interacción de las nanopartículas con las ondas electromagnéticas es lo que permite el apantallamiento, mientras que el polímero se encarga del manejo estructural y contribuye indirectamente en la absorción o el reflejo de la interferencia electromagnética [1], [12].

En este trabajo se realizó un análisis sobre nanocompuestos SWCNT/poliimida, con el objetivo de dar una interpretación física a la relación del apantallamiento electromagnético de diferentes muestras medido a partir de parámetros S, por medio del uso de descriptores morfológicos propuestos y medibles sobre los datos de las mismas. Estos datos son obtenidos a partir de microscopía de Fuerza de Sonda Kelvin de Segundo Armónico (KPFM) [5], capaz de realizar la medición del gradiente de capacitancia $\partial C/\partial z$ en cada posición. Son representaciones de la muestra en formato de imagen y sobre las mismas se realiza la medición de descriptores. La búsqueda de relaciones entre los descriptores y la medición de efectividad de apantallamiento (SE) se realiza usando modelos de Machine Learning (ML), capaces de detectar las características más influyentes en la propiedad de interés a partir de modelos predictivos e interpretables.

Con base en lo expuesto, la pregunta de investigación que dirige este proyecto es: ¿Cuáles son los descriptores morfológicos de la red de SWCNTs, extraídos de imágenes KPFM, que se relacionan cuantitativamente con la eficacia de apantallamiento electromagnético del nanocompuesto? La hipótesis es que la morfología de la red, caracterizada a partir de las imágenes de gradiente de capacitancia, determina en parte significativa el SE del material, y que modelos de ML interpretables pueden identificar y jerarquizar esa relación.

Si bien no es nueva la caracterización de muestras a partir de descriptores apuntando a la relación con SE, realizada en el trabajo de Castañeda Uribe y Avila [4], usando tres descriptores morfológicos (CNTD, CNTA, CNTS) y parámetros S, la conexión sistemática con la morfología de la red no está completamente establecida en este tipo de nanocompuestos. 

Este trabajo extiende esta línea a partir de tres enfoques:(i) un conjunto ampliado de descriptores morfológicos, soportado por un pipeline de procesamiento de imágenes diseñado para extraerlos de forma robusta y generalizada; (ii) un espacio experimental con mayor variabilidad en las condiciones de fabricación; y (iii) la integración de modelos de ML interpretables para identificar y jerarquizar los descriptores que se relacionan con el SE, habilitando su posterior interpretación física. 

El enfoque del trabajo busca evitar caracterizaciones cualitativas, fundamentando el análisis físico a proponer. Se reconocen limitaciones en cuanto a la cantidad de muestras y a la proyección 2D de la red tridimensional, derivada del procesamiento de imágenes. De resultar efectivo, se podría generalizar este proceso para el análisis de muestras de nanocompuestos poliméricos similares.

\section{Objetivos}
El objetivo general de este trabajo es analizar la relación entre descriptores morfológicos obtenidos a partir de imágenes de distribuciones de nanotubos de carbono (CNTs) y la eficacia de apantallamiento electromagnético en nanocompuestos poliméricos de CNT/poliimida, mediante el uso de modelos de inteligencia artificial interpretables y mediciones experimentales de parámetros S. Para ello se plantean los siguientes objetivos específicos:
\begin{itemize}

\item Aplicar técnicas de procesamiento digital de imágenes para reducir el ruido de imágenes obtenidas por KPFM, con el fin de mejorar el contraste entre los nanotubos de carbono y la matriz polimérica.
\item Extraer descriptores morfológicos físicamente interpretables a partir de las imágenes filtradas, que representen cuantitativamente la distribución espacial de los CNTs en la matriz.
\item Entrenar modelos de inteligencia artificial explicables utilizando los descriptores morfológicos como variables de entrada y las eficacias de apantallamiento electromagnético, obtenidas mediante parámetros S, como variable objetivo.
\item Interpretar la influencia de la morfología interna de los nanocompuestos en su eficacia de apantallamiento electromagnético, empleando descriptores morfológicos extraídos de imágenes de caracterización para explicar el comportamiento del material desde una perspectiva fundamentada en la física del electromagnetismo y de los semiconductores.
\end{itemize}
